% Generated draft: requires mathematical and source review.
\documentclass[11pt]{article}
\usepackage[margin=25mm]{geometry}
\usepackage{fontspec}
\setmainfont{Latin Modern Roman}
\usepackage{amsmath,amssymb,mathtools,unicode-math}
\setmathfont{Latin Modern Math}
\usepackage{enumitem,needspace,xurl,hyperref}
\setlength{\parindent}{0pt}
\setlength{\parskip}{5pt}
\setlength{\emergencystretch}{3em}
\newcommand{\N}{\mathbb{N}}
\newcommand{\Z}{\mathbb{Z}}
\newcommand{\Q}{\mathbb{Q}}
\newcommand{\R}{\mathbb{R}}
\newcommand{\C}{\mathbb{C}}
\newcommand{\F}{\mathbb{F}}
\newcommand{\A}{\mathbb{A}}
\newcommand{\PP}{\mathbb{P}}
\newcommand{\E}{\mathbb{E}}
\newcommand{\eps}{\varepsilon}
\newcommand{\dd}{\,\mathrm{d}}
\begin{document}
\section*{UnsolvedMath ID: 3000012}
\textbf{Source problem code:} catalogue entry
\subsection{Definitions and mathematical statement}
Let $H=(V,E)$ be a hypergraph, with edges regarded as subsets of $V$. A matching is a family $\mathcal D\subseteq E$ of pairwise disjoint edges. A vertex cover is a set $C\subseteq V$ satisfying $C\cap e\ne\varnothing$ for every $e\in E$.

The K\H{o}nig property means that there exist a matching $\mathcal D$ and a choice $c(e)\in e$ for each $e\in\mathcal D$ such that
\[
C=\{c(e):e\in\mathcal D\}
\]
is a vertex cover. Thus the cover consists exactly of one selected vertex from each matching edge, with no additional vertices; equality of infinite cardinalities alone is not the definition.

\textbf{Compactness conjecture (R. Aharoni and N. Bowler, independently).} Suppose every edge of $H$ is finite and, for every finite subfamily $E'\subseteq E$, the hypergraph $(V,E')$ has the K\H{o}nig property. Then $H$ has the K\H{o}nig property.

The finite-subfamily condition retains the entire vertex set $V$. Its witnessing matching and cover may depend on $E'$ and need not come with any prescribed compatibility. No bound on edge sizes, vertex degrees, or the cardinalities of $V$ and $E$ is imposed. Empty edges are automatically excluded by the hypothesis, since a singleton subfamily containing one could have no vertex cover. The source recalls finite bipartite graphs as the classical K\H{o}nig-theorem case and reports no verified resolution of this infinite-hypergraph conjecture.
\subsection{Short English statement}
If every finite selection of finite hyperedges has a matching with a one-vertex-per-matched-edge cover, must the entire hypergraph have such a matching and cover? The conjecture allows arbitrary infinite hypergraphs and unbounded finite edge sizes.
\subsection{Sources}
\begin{itemize}
\item Egres Open, \emph{Compactness of K\H{o}nig-property}, listed at \url{https://oldlemon.cs.elte.hu/egres/open/List_of_open_problems}; HTML extraction, accessed 2026-07-29. The record warns that the archived open-problem index is not a current resolution authority.
\item R. Aharoni and N. Bowler, independent conjecture attribution in the source. The statement references p. 19, Problem 6.7, identified by its research summary as Aharoni's book; no title, year, or edition is supplied.
\item Local registry \texttt{data/unsolvedmath/problems.json}, permanent ID $3000012$, source code catalogue entry, \emph{Compactness of K\H{o}nig-property}. The record labels this open and OPEN-TRIAGE, with NEEDS\_REVIEW status and rights notes; the archived footer did not display a redistribution license according to the source record.
\end{itemize}
\end{document}
